%% ============================================================
\section{Experiments}
\label{sec:experiments}
%% ============================================================

We ask three questions. First, does roadmap-tilted sampling plan better than
standard trajectory optimisers on the same world model and the same budget
(\S\ref{sec:exp:planning})? Second, does any advantage survive as the planning
horizon grows (\S\ref{sec:exp:horizon})? Third, does the roadmap cost-to-go
contribute anything beyond the $L_2$ planning objective that the optimiser is
already descending (\S\ref{sec:exp:ablation})?

\subsection{Setup}
\label{sec:exp:setup}

\paragraph{Environments and world models.}
We evaluate on five environments spanning two world-model families. Reacher and
Cube use LeWM/JEPA latents, in which each frame is encoded as a single $192$-d
token. Point maze, Push-T and Wall use DINO-WM, where a frame becomes $196$ ViT
patch tokens. The world model is frozen throughout; no component of any planner
is trained.

\paragraph{Protocol.}
Unless stated otherwise every planner receives a matched budget of $150$
optimiser iterations, and CEM is additionally tuned to $150$ world-model
rollouts ($30$ samples $\times$ $5$ steps) so that it matches gradient descent's
rollout count rather than its stock $9000$. Each cell reports $100$ seeds, one
episode per seed. Goals are drawn per environment: \texttt{random\_state} on
Reacher, Point maze and Wall; from the dataset on Push-T, whose random
start/goal pairs are frequently unconnectable and floor every planner near
zero. The goal protocol therefore is not uniform across environments and we
report it per row.

\paragraph{Horizons.}
$\mathrm{goal}_H$ counts world-model steps, each covering \texttt{frameskip}$=5$
environment steps, so $H\in\{5,10,15\}$ corresponds to $25$, $50$ and $75$
environment steps. We note in \S\ref{sec:exp:limits} that this makes $H{=}15$
longer than the episode cap on two environments.

\paragraph{Particle budgets are not uniform, and this favours the baselines.}
On the LeWM environments our sampler runs $8$ particles with $8$ smoothing
samples. On the DINO-WM environments it runs $3$ particles with smoothing
disabled, and $2$ at $H{=}15$. This is a hardware limit rather than a tuning
choice: DINO-WM retains $196$ patch tokens per frame through the rollout for
autograd, and we measured that $4$ particles exceeds $16$\,GB of VRAM at
$H{=}10$ and $3$ particles at $H{=}15$. Our method is therefore
\emph{reduced-budget} exactly where it is compared against baselines that are
not.

\subsection{Does it plan better?}
\label{sec:exp:planning}

Table~\ref{tab:planners} reports success rates across the shared grid. Our
method is strongest on Reacher, where it reaches $100\%$, $92\%$ and $50\%$ at
$H=5,10,15$ against best-baseline values of $62\%$, $38\%$ and $32\%$ --- a
margin of $38$, $54$ and $18$ points. On Point maze it leads at the two shorter
horizons ($75$ vs $54$; $81$ vs $67$) and \emph{loses} at $H{=}15$ ($61$ vs CEM's
$69$). On Push-T it leads the single available baseline at $H{=}5$ ($68$ vs
$57$).

Table~\ref{tab:cube} reports OGBench Cube under the evaluation protocol of its
reference implementation: goals drawn $25$ environment steps ahead in an expert
episode, $50$ seeds, and success when the block lies within $0.04$\,m of its
target. Our method leads on both views, $88\%$ over all seeds and $79\%$ on the
non-trivial subset, against $82/69$ for gradient descent and $76/59$ for the
reference CEM solver. Paired per seed on the non-trivial subset it wins $3$,
loses $0$ and ties $26$ against gradient descent: it solves a strict superset,
rather than trading different seeds. We emphasise two caveats that make this a
harder comparison than the headline suggests. The gradient-descent baseline is
our own tuned configuration --- the generic one \emph{diverges} on this
environment, moving away from the goal --- so the comparison is against a
strengthened rather than a default baseline. And $42\%$ of goals under this
protocol are trivial: the block already lies inside the success radius, so the
goal is satisfied before any planning and every planner scores them. This
follows from the reference protocol, which samples start indices uniformly with
no filter on block motion, and it is why we report the non-trivial column.

\subsection{Does it generalise across horizons?}
\label{sec:exp:horizon}

Only partially, and the pattern differs by environment. On Reacher the advantage
narrows monotonically with horizon --- $+38$, $+54$, $+18$ points over the best
baseline --- and absolute performance falls from $100\%$ to $50\%$. On Push-T it
collapses outright, $68\%\to27\%\to5\%$. On Point maze it is roughly flat
($75\to81\to61$) but the ranking inverts: we lead at $H\le10$ and trail CEM at
$H{=}15$. Across the grid, CEM is the planner whose performance is most stable
in the horizon, and on Wall it \emph{improves} with horizon ($33\to40\to49$)
while GRASP degrades monotonically ($21\to9\to5$).

We do not claim horizon-robustness. Two confounds prevent a clean reading, and
both act against us. Our particle budget \emph{shrinks} with horizon on the
DINO-WM environments ($3$ particles at $H\le10$, $2$ at $H{=}15$) for the memory
reasons above, so the longest-horizon cells are also the ones where our sampler
has the least population to select among. And on Reacher and Cube, $H{=}15$
corresponds to $75$ environment steps against a $50$-step episode cap, so those
cells plan past the end of the episode. The honest summary is that the method's
advantage is clearest at short and medium horizons and is not established at
long ones.

\subsection{Does the roadmap cost-to-go add anything over the $L_2$ objective?}
\label{sec:exp:ablation}

The tilt potential and the optimiser's objective are different functions of the
same rollout, so it is reasonable to ask whether the roadmap solve is doing any
work that the $L_2$ planning loss could not. We test this directly by replacing
the cost-to-go $V$ in the importance weights with the $L_2$ objective the
Langevin kernel already descends, leaving every other component untouched.

The comparison is run on Point maze at $H{=}10$ over the $38$ seeds where
gradient descent fails and our method succeeds. These seeds are chosen because
gradient descent \emph{is} the $L_2$ objective, so the objective alone provably
does not solve any of them; our method solves all $38$, giving full headroom to
fall.

Replacing $V$ with the $L_2$ objective loses $9$ of $38$ seeds ($38/38 \to
29/38$). Counting failures at the iteration cap, it also converges far more
slowly: mean $25.6$ versus $47.6$ iterations to first success, median $5$ versus
$11$, winning $19$ seeds and losing $5$. We report the censored statistic
because conditioning on seeds both arms solve discards precisely the hard cases
the ablation is about.

The diagnostics identify the mechanism. The roadmap tilt sustains roughly twice
the selection pressure of the objective tilt (cumulative selection KL $0.73$ vs
$0.36$ nats), resamples more often ($1.9$ vs $1.3$), and drives the effective
sample size \emph{lower} ($1.84$ vs $2.27$ of $3$), i.e.\ it actually
concentrates mass, whereas the objective leaves the population near-uniform. On
two seeds the objective tilt is completely inert; the roadmap tilt never is.
The explanation is structural: the kernel already descends the $L_2$ objective,
so optimised particles have similar objective values and little spread remains
to tilt on, while $V$ is a different function of the state and still separates
particles the objective can no longer distinguish. Consistently, the annealing
exponent $\beta$ runs \emph{lower} for the roadmap tilt ($0.897$ vs $0.970$):
the objective arm pushes its exponent higher to compensate and still achieves
half the selection pressure.

\paragraph{A negative result.}
This contribution does not transfer to Cube. There the correlation between $V$
and latent distance is $0.003$, the tilt is near-inert (final ESS $28.6/32$,
cumulative KL $0.14$), and a four-arm grid confirms the diagnosis: raising
$\beta_{\max}$ does activate the tilt (KL $0.71$, ESS $19.2$) and changes
nothing, with $4/5$ non-trivial seeds solved in every arm and the same seed
failing throughout; a four-fold denser roadmap leaves the off-manifold distance
unmoved. Our method still leads on Cube (\S\ref{sec:exp:planning}), but it does
so through the gradient kernel rather than the roadmap, and we do not claim
otherwise. We conjecture that the LeWM Cube latents encode the arm
configuration but not the block position --- the quantity the task and its
success criterion depend on --- which would make a roadmap built in that latent
space unable to represent the relevant distance.

\subsection{Limitations}
\label{sec:exp:limits}

\textbf{Missing baselines.} Eight of the nine Push-T baseline cells and all
Wall cells for our method are unscheduled, so the Push-T results at $H\ge10$ and
the Wall column stand without a complete comparison.
\textbf{Horizon versus episode length.} On Reacher and Cube, $H{=}15$ plans $75$
environment steps against a $50$-step cap; those cells should be read with that
in mind.
\textbf{Run-to-run variance.} The planner is not reproducible at fixed seed:
\texttt{seed} fixes the start/goal draw but not the particle initialisation,
Langevin noise, or the cost-to-go surrogate fit. Re-running a completed arm
reproduced $38/38$ outcomes exactly, but on marginal seeds we measured
disagreement of roughly $3$--$4$ points per $100$-seed cell. Per-seed claims
should be read accordingly.
\textbf{Trivial goals on Cube.} $42\%$ of goals are satisfied at $t{=}0$; the
non-trivial column is the discriminating one.
